UART is used as the reporting interface between the FPGA-side verification platform and the host computer. After packets are injected, checked, and classified on chip, the generated verification reports are sent through the UART interface. In this verification flow, UART is not part of the NoC data path; it is only used to transfer test results and status information to the PC.

The reported information includes packet-level and case-level events. Packet-level reports indicate whether an injected packet has been successfully received or failed due to an error condition. Case-level reports summarize the execution result of each test case, including pass/fail information and latency-related statistics. Progress and completion messages can also be returned to indicate the current execution state of the verification platform.

On the PC side, the received UART messages are recorded and parsed into structured result files. These results can then be used to analyze packet correctness, error distribution, and latency behavior. By moving result collection and analysis to the PC side, the FPGA logic only needs to perform essential checking and reporting, while more detailed statistical analysis can be handled after execution.
